\newpage
\phantomsection
\label{prf:boolean-ring-characteristic-2}

\begin{remark*}[Return]
\hyperref[def:boolean-ring]{Return to Definition: Boolean Ring}
\end{remark*}

\begin{remark*}[Relationship to existing proof]
The proof of this theorem is contained in full in
\hyperref[prf:boolean-ring-self-additive-inverse]{Self Additive Inverse
(prf:boolean-ring-self-additive-inverse)}, which proves $p + p = 0$ directly
from idempotence. The label \texttt{prf:boolean-ring-characteristic-2} is the
entry point linked from the Boolean Ring definition; the content is shared.
\end{remark*}

\begin{theorem*}[Every Boolean Ring Has Characteristic~2]
In any Boolean ring $R$,
\[
p + p = 0 \qquad \text{for all } p \in R.
\]
Equivalently, $\operatorname{Characteristic2}(R,+,0)$ holds in every
Boolean ring.
\end{theorem*}

\begin{proof}
\textbf{Professional Standard Proof.}~\\
Apply idempotence to the element $p + p$:
\[
(p + p) \cdot (p + p) = p + p. \tag{i}
\]
Expand the left side using right distributivity~(9), then left
distributivity~(8), then collapse each $p \cdot p = p$ by
idempotence~(11):
\[
(p + p) \cdot (p + p)
  \;\overset{(9)}{=}\; p\cdot(p+p) + p\cdot(p+p)
  \;\overset{(8)}{=}\; (p^2 + p^2) + (p^2 + p^2)
  \;\overset{p^2=p}{=}\; (p + p) + (p + p). \tag{ii}
\]
Combining (i) and (ii): $(p+p) + (p+p) = p+p$.
Setting $s = p+p$, the equation $s + s = s$ gives $s = 0$
by adding $-s$ and applying associativity, the inverse law,
and the identity law.
Hence $p + p = 0$.
\end{proof}

\begin{remark*}[Dependencies]~\\
\begin{itemize}
  \item \hyperref[def:boolean-ring]{Definition: Boolean Ring}
  \item \hyperref[def:idempotent-element]{Definition: Idempotent Element}
\end{itemize}
\end{remark*}
